\section{Additional Training and Execution Studies}
\label{app:additional}

\subsection{Paired Training-Recipe Study}
The recipe study tests two development combinations, G (frozen DINOv3-S)
and I (trainable CNN), both with P07 supervision, the full visual window,
and current-target RGB reconstruction. For each profile, training seeds
0 and 1 compare global batch 16 on sixteen GPUs against batch 64 on
thirty-two GPUs. The larger-batch recipe doubles the learning rate from
$2\!\times\!10^{-4}$ to $4\!\times\!10^{-4}$ and reduces warm-up from 1,000
to 250 updates, preserving 16,000 warm-up samples. This is a joint recipe
and resource comparison, not an isolated batch-size intervention.

All eight runs complete four epochs and 641,664 samples. B16 uses 40,104
updates; B64 uses 10,028, retaining each epoch's 32-sample tail without
padding, repetition, or dropping. AdamW, clipping, FP32/TF32, losses,
normalization, and the four-epoch cosine exposure are held fixed.
All final checkpoints undergo the same sixteen-device full350 validation
and the original 1,400-scene evaluation with paired scene and inference seeds.
The seed-0 B16 runs are fresh reruns: their scene-level success outcomes
exactly reproduce the corresponding BASE2 endpoints. They do not add a new
independent training seed. Seed 1 supplies the additional realization.

% Generated from Research BATCH0/results/full_measurements_20261016.json.
\begin{table*}[t]
\centering\small
\caption{\tablelang{Completed training-recipe study, four epochs per run. G uses frozen
DINOv3-S and I a trainable CNN; both use the corresponding P07 combination.
Batch 64 jointly changes learning rate and device count (details in text).
Training time excludes saving and validation. Job time runs from job creation
through complete evaluation, including the observed queue and startup.}{已完成的训练配方研究，每次训练四轮。G 使用冻结的 DINOv3-S，I 使用可训练 CNN；二者均采用对应的 P07 组合。Batch 64 同时改变学习率和设备数（详见正文）。训练耗时不含保存和验证。作业耗时从创建作业起计至完整评测结束，包含实际发生的排队和启动时间。}}
\label{tab:batch_recipe}
\begin{tabular}{@{}lrrrrrrr@{}}
\toprule
\tablelang{Profile}{配置} & \tablelang{Seed}{种子} & \tablelang{Batch}{批大小} & \tablelang{GPUs}{GPU 数} & \tablelang{Success /1,400}{成功数 /1,400} & \tablelang{SR (\%)}{成功率 (\%)} & \tablelang{Train (min)}{训练 (min)} & \tablelang{Job (min)}{作业 (min)} \\
\midrule
G & 0 & 16 & 16 & 849 & 60.64 & 263.40 & 301.28 \\
G & 0 & 64 & 32 & 855 & 61.07 & 99.27 & 121.32 \\
G & 1 & 16 & 16 & 834 & 59.57 & 263.28 & 303.58 \\
G & 1 & 64 & 32 & 822 & 58.71 & 97.48 & 120.23 \\
I & 0 & 16 & 16 & 935 & 66.79 & 220.75 & 257.20 \\
I & 0 & 64 & 32 & 884 & 63.14 & 78.86 & 101.74 \\
I & 1 & 16 & 16 & 857 & 61.21 & 224.91 & 258.42 \\
I & 1 & 64 & 32 & 935 & 66.79 & 79.80 & 101.02 \\
\bottomrule
\end{tabular}
\end{table*}

% Generated from the 11,200 original BATCH0 episode results.
\begin{table*}[t]
\centering\small
\caption{\tablelang{Training-recipe successes per task (100 scenes per entry).
B16 and B64 include their respective learning-rate/device settings;
s0/s1 denote training seeds. All runs use the same accepted scene identities
and inference seeds. Differences across tasks remain visible despite similar totals.}{训练配方研究的逐任务成功数（每项 100 个场景）。B16 和 B64 分别包含各自的学习率与设备设置；s0/s1 表示训练 seed。所有运行使用相同的已接受场景身份和推理 seed。总分相近时，逐任务差异仍然存在。}}
\label{tab:batch_tasks}
\begin{tabular}{@{}lrrrrrrrr@{}}
\toprule
 & \multicolumn{4}{c}{\tablelang{G: frozen DINOv3-S}{G：冻结 DINOv3-S}} & \multicolumn{4}{c}{\tablelang{I: trainable CNN}{I：可训练 CNN}} \\
\cmidrule(lr){2-5}\cmidrule(l){6-9}
\tablelang{Task}{任务} & B16/s0 & B64/s0 & B16/s1 & B64/s1 & B16/s0 & B64/s0 & B16/s1 & B64/s1 \\
\midrule
\tablelang{Grasp}{抓取} & 92 & 87 & 91 & 89 & 90 & 95 & 87 & 96 \\
\tablelang{Belt}{传送带} & 100 & 100 & 100 & 97 & 97 & 100 & 99 & 100 \\
\tablelang{Occlusion}{遮挡} & 100 & 99 & 100 & 100 & 97 & 98 & 97 & 100 \\
SO(3) & 87 & 88 & 68 & 53 & 86 & 75 & 56 & 55 \\
XYZ & 75 & 58 & 60 & 69 & 82 & 69 & 77 & 85 \\
\tablelang{SE(3) placement}{SE(3) 放置} & 0 & 0 & 0 & 0 & 0 & 0 & 1 & 0 \\
\tablelang{Trans. placement}{平移放置} & 1 & 10 & 10 & 4 & 4 & 5 & 13 & 29 \\
\tablelang{Full passage}{完整穿越} & 0 & 0 & 0 & 0 & 0 & 0 & 0 & 0 \\
\tablelang{Phase passage}{相位穿越} & 0 & 0 & 13 & 0 & 72 & 32 & 14 & 46 \\
\tablelang{Can}{罐体} & 99 & 99 & 97 & 98 & 100 & 98 & 99 & 100 \\
\tablelang{Slider}{滑块} & 83 & 81 & 76 & 80 & 81 & 86 & 75 & 89 \\
\tablelang{Sphere}{球体} & 40 & 50 & 50 & 57 & 56 & 47 & 60 & 73 \\
\tablelang{Tabletop sphere}{桌面球体} & 100 & 99 & 99 & 100 & 98 & 99 & 99 & 100 \\
\tablelang{Replan}{重规划} & 72 & 84 & 70 & 75 & 72 & 80 & 80 & 62 \\
\midrule
\tablelang{Total}{总计} & 849 & 855 & 834 & 822 & 935 & 884 & 857 & 935 \\
\bottomrule
\end{tabular}
\end{table*}


\paragraph{Performance and variability.}
For G, B64 minus B16 gives $+0.43$ and $-0.86$ percentage points at seeds
0 and 1, with paired 95\% intervals $[-1.21,2.07]$ and $[-2.50,0.79]$.
For I, the changes reverse direction: $-3.64$ $[-5.64,-1.64]$ and
$+5.57$ $[3.57,7.57]$ points. At B16, I falls from 935 to 857 successes
between seeds; phase passage falls from 72 to 14. Neither the original
seed-0 score nor the mean response to a recipe change establishes stability.
The matched E/F controls have not been repeated at seed 1.

The two-seed mean recipe differences are $-0.21$ $[-1.36,0.96]$ points
for G and $+0.96$ $[-0.50,2.39]$ for I. Intervals use 20,000 percentile
bootstrap samples, RNG seed 20261017, resampling 100 scene identities within
each task and retaining their batch and training-seed pairing. They condition
on these two trained seeds and fourteen tasks; they do not estimate variation
over new training seeds. Table~\ref{tab:batch_tasks} also shows task losses
hidden by aggregate scores: G loses 17 XYZ successes at seed 0 and 15 SO(3)
successes at seed 1. Auxiliary metrics are not uniformly preserved; for
example, I's pose rotation error increases by 6.17\% and 14.35\%.

\paragraph{Time and adoption.}
Mean job-creation-to-result time decreases from 302.43 to 120.77 minutes
for G and from 257.81 to 101.38 minutes for I. These measured times include
queueing, setup, training, saving, validation, export, and evaluation.
They do not identify a hardware-only speedup. Training-compute GPU-hours
decrease from 70.22 to 52.47 for G and from 59.42 to 42.31 for I;
this measure excludes saving and validation and is not a billing record.
The predeclared working tolerance was a three-point aggregate SR loss.
G satisfies it for both observed seeds, whereas I does not. The project
therefore retains B16 as the general recipe; B64 is a conditional option
for this G configuration when aggregate SR is the criterion. It is not
adopted for I, task-sensitive comparisons, or an untested final design.

\subsection{Frozen Three-State Gripper Readout}
\label{app:learned_decoder}
The execution study freezes the fifth-epoch P03 policy and learns a
three-state readout. Its 889-dimensional input concatenates the current
action hidden state at the final Euler evaluation ($\tau=0.8$), eight
frames of available end-effector pose, gripper opening and validity,
the sampled six-dimensional arm-action horizon, and the current raw gripper
value. A 128-hidden-unit GELU MLP predicts open/hold/close logits; argmax
replaces only the executed slot-0 gripper command. A second head with the
same hidden width reads only the raw gripper value. No contact truth or
future label is added to the observation. The frozen policy, normalization,
five Euler steps, and synchronous single-step execution remain common.

\paragraph{Supervision and offline evaluation.}
The original 11,200 training and 350 validation demonstrations supply
1,324,374 and 41,533 valid issued-command labels; training and validation
source seeds are disjoint within each task. Every valid current command is retained, including windows whose
later action slots are invalid. Only capture demonstrations contain hold
labels: 2,562 in training (0.1934\% of all labels) and 79 in validation.
The other thirteen tasks have no valid hold labels; all phase-passage
training and validation labels are close. This limits what the decoder
can learn about useful holds outside capture.

Both heads use seed 0, twenty epochs, AdamW with learning rate
$3\!\times\!10^{-4}$ and weight decay $10^{-4}$, batch 4,096, gradient
clipping at 1, and ordinary cross-entropy without class reweighting.
Feature means and variances use training data only. Each head is selected
by the lowest validation mean of the three per-class NLLs, choosing the
earliest epoch in a tie. Context selects epoch 15 with NLL 0.01330;
raw-only selects epoch 20 with NLL 0.19640.

% Generated from Research GRIPPER2/results/learned/training-report.json.
\begin{table}[t]
\centering\small
\caption{Frozen-decoder validation confusion matrices: 41,533 valid issued
commands from 350 demonstrations. Rows are ground truth; columns are predicted
open/hold/close. These are offline classification counts, not policy successes.}
\label{tab:decoder_confusion}
\begin{tabular}{@{}lrrrrrr@{}}
\toprule
 & \multicolumn{3}{c}{Raw only} & \multicolumn{3}{c}{Context} \\
\cmidrule(lr){2-4}\cmidrule(l){5-7}
GT & Open & Hold & Close & Open & Hold & Close \\
\midrule
Open & 8003 & 1 & 214 & 8132 & 0 & 86 \\
Hold & 0 & 75 & 4 & 0 & 79 & 0 \\
Close & 135 & 9 & 33092 & 105 & 4 & 33127 \\
\bottomrule
\end{tabular}
\end{table}

The context head recalls 79/79 hold labels with precision 79/83, versus
75/79 and 75/85 for raw-only (Table~\ref{tab:decoder_confusion}). It correctly
classifies 137/153 one- or two-tick pulses, compared with 93/153; neither
head incorrectly maps these pulses to hold. These counts describe frozen
features on demonstration observations, not closed-loop policy performance.
All three per-class NLLs improve. The 79 hold labels receive one third of
the checkpoint-selection metric and account for 74.5\% of its reduction;
training cross-entropy remains unweighted. The context head still makes
86 open-to-close and 105 close-to-open validation errors, so the mean loss
does not constrain the cost of an individual closed-loop decision.

\paragraph{Predefined development cases.}
The context head is tested on eleven known development cases. It succeeds
on E5 capture scenes 0/9/14/80/85, E5 replan scenes 49/88, and E7 XYZ scene
64 and replan scene 7; the E7 cases reuse the E5-trained head without
further fitting. It fails on E5 capture scene 52 and replan scene 35.
The head succeeds in 9/11 cases, below the predefined all-eleven
criterion, and is therefore excluded from the final policy. The case count
describes this selected development set.

An observation-only control preserves the original commands and random
state. Twenty-eight paired sign/dead-zone trials reproduce their reference
outcomes. Offline recomputation of 608 issued commands gives
a maximum CPU-logit error of $3.82\!\times\!10^{-6}$ and argmax commands
agree exactly. The other six action dimensions and predicted horizons also
pass the recorded comparisons. Complete per-frame policy-input tensors
were not saved, so this is not a full input-tensor equality check.

\begin{figure*}[t]
\centering
\includegraphics[width=\textwidth]{decoder_failures.pdf}
\caption{Frozen gripper-readout failures on the two failed predefined
development cases. Commands and measured opening use physical time; event
strips retain each run's contact and terminal outcome. Sign, the fixed
0.1 dead zone, and the learned context head use the same frozen policy.
The head's first intervention reverses opening into closing, before the
observed physical divergence. Detailed windows and full-run events serve
different time ranges, labeled on their axes. Contact refers to actual
bilateral collision; replan attachment is a separate event. Raw predictions
are sampled at 20\,Hz and physical states every 2\,ms. Separation of the learned
run from Sign uses a 10\,$\mu$m aperture or componentwise position threshold. The cases are
selected development diagnostics, not new-scene success estimates or
future-prediction interventions.}
\label{fig:decoder_failures}
\end{figure*}

% Balance the final text page without changing the full-page figure size.
\newpage
\paragraph{Physical failure sequence.}
Figure~\ref{fig:decoder_failures} shows both failed cases. In each, the first
intervention reverses an approximately $-1$ open command into $+1$ close;
it is not merely suppression of a near-zero output. Capture changes at
0.650\,s, followed by opening divergence at 0.666\,s and no bilateral contact.
Replan applies the altered command at 1.550\,s, after its native one-tick
queue; opening diverges at 1.556\,s. The learned head fails at 1.324\,s
(tray contact) and 5.174\,s (downstream exit), respectively. Sign succeeds
in both; the fixed dead zone succeeds in replan but also fails in capture.
Recorded predictions, commands, and 2\,ms physical fields match their
references before intervention. Subsequent arm and object divergence
reflects closed-loop feedback. These records establish failures of this
readout under the recorded conditions; they do not test a future-to-action
prediction interface or establish decoder performance on unseen scenes.

\paragraph{Retained-feature diagnostics.}
At the first reversal in capture 52 and replan 35, replacing only the
action hidden state with the preceding frame's value changes the
close-minus-open logit margin from $+2.452/+0.700$ to $-13.839/-9.811$.
Conversely, inserting the current hidden state into the preceding feature
vector yields $+2.243/+3.184$; replacing only the raw value does not remove
either reversal. These substitutions measure a local response of the
frozen classifier to mixed features, which need not correspond to
realizable observations or successful executions. The first reversal
precedes any prior intervention by this head. The retained data do not
separate demonstration-to-policy feature mismatch from sampling
sensitivity of the final Euler hidden state; both explanations remain
untested. The two failures remain unexplained by a single isolated mechanism.
